\documentclass[11pt]{article}
\usepackage[margin=1in]{geometry}
\usepackage{amsmath,amssymb,amsthm,mathtools,booktabs,longtable,array}
\usepackage[T1]{fontenc}
\usepackage{lmodern}
\usepackage[hidelinks]{hyperref}
\usepackage{microtype}
\newtheorem{theorem}{Theorem}
\newtheorem{lemma}[theorem]{Lemma}
\newtheorem{corollary}[theorem]{Corollary}
\newcommand{\spn}{\operatorname{spread}}
\newcommand{\rank}{\operatorname{rank}}
\newcommand{\tr}{\operatorname{tr}}
\newcommand{\diag}{\operatorname{diag}}
\newcommand{\R}{\mathbb R}
\newcommand{\J}{J}
\title{SP-08: Centered Rank-Two Sign Matrices}
\author{Sidney Holden\\{\small Center for Computational Biology, Flatiron Institute}\\{\small Simons Foundation}}
\date{13 September 2026}
\hypersetup{pdfauthor={Sidney Holden}}
\begin{document}
\maketitle
\begin{abstract}
This continuation proves the conjectured rank-two candidate bound for every
endpoint matrix whose centered sign matrix has rank at most two, in every
dimension and throughout the normalized interval $-1\le a<1$. An analytic
monotonicity lemma retains integer block sizes, avoiding a fractional-mass
relaxation. The negative-parameter comparison is analytic; the nonnegative
comparison is certified by exact rational polynomial identities with an
independent standard-library verifier. A consequence is the unrestricted strip
$-1\le a\le-1+(1+\sqrt2)/n^2$, improving the coefficient $1/3$ in the supplied
boundary-strip argument. \textbf{The unrestricted SP-08 conjecture is not
proved here.} Higher-rank indefinite deficits remain outside the argument.
\end{abstract}

\section{Problem, candidate, and status}
Let $\mathcal S_n[a,1]$ be the real symmetric $n\times n$ matrices all of whose
entries, including diagonal entries, belong to $[a,1]$. Here $n\ge2$ and
$-1\le a<1$. For a symmetric matrix,
\[
 \spn(A)=\lambda_{\max}(A)-\lambda_{\min}(A).
\]
SP-08 asks whether an endpoint matrix of rank two always attains the maximum
spread over this box~\cite{sp08}. Convexity gives an endpoint maximizer, but does
not give its rank.

Put $d=1-a$. If $z$ indicates a $k$-element subset, the matrix
$B_k=J-dzz^T$, $1\le k<n$, has two nonzero eigenvalues of opposite signs.
Its squared spread is
\begin{align}
 F(k)&=n^2+d\{2nk-(a+3)k^2\},\label{eq:F}\\
 M_n(a)^2&=\max_{1\le k<n,\ k\in\mathbb Z} F(k).
\end{align}
Completing the square gives
\begin{equation}\label{eq:round}
 F(t)=\frac{4n^2}{a+3}-d(a+3)\left(t-\frac{n}{a+3}\right)^2.
\end{equation}
An allowed integer nearest $n/(a+3)$ is optimal in this family. In particular,
\begin{equation}\label{eq:U}
 M_n(a)^2\ge U:=\frac{4n^2}{a+3}-\frac{d(a+3)}4.
\end{equation}
The term subtracted in \eqref{eq:U} is essential: the continuous optimum alone
would not establish the finite-dimensional conjecture.

\begin{theorem}[Centered sign rank at most two]\label{thm:centered}
Set $c=(1+a)/2$ and $r=(1-a)/2$. If $A\in\mathcal S_n[a,1]$ is an endpoint
matrix and its sign matrix $S=(A-cJ)/r$ has rank at most two, then
$\spn(A)\le M_n(a)$.
\end{theorem}

\begin{theorem}[Unrestricted boundary strip]\label{thm:strip}
For every $n\ge2$ and
\[
 -1\le a\le -1+\frac{1+\sqrt2}{n^2},
\]
$B_{\lfloor n/2\rfloor}$ maximizes spread over all of $\mathcal S_n[a,1]$.
There is no rank or inertia restriction in this statement.
\end{theorem}
Theorem~\ref{thm:strip} uses only the analytic, negative-parameter portion of
Theorem~\ref{thm:centered}; it does not depend on a numerical test or the
nonnegative-parameter certificate.

\section{Classification and compression}
A symmetric rank-two sign matrix is indefinite and has both positive and
negative diagonal entries. Indeed, if all diagonal signs were equal, every
principal $2\times2$ minor would vanish; their sum, the product of the two
nonzero eigenvalues, cannot vanish. Choose a positive and a negative diagonal
position and switch one sign so that their principal matrix is
\[
 H_2=\begin{pmatrix}1&1\\1&-1\end{pmatrix},\qquad H_2^{-1}=H_2/2.
\]
The vanishing Schur complement classifies all columns by their two entries
on these positions: they are $\pm(1,1)^T$ or $\pm(1,-1)^T$. Consequently,
up to permutation,
\begin{equation}\label{eq:class}
 S=\Delta\begin{pmatrix}J_p&J_{p,q}\\J_{q,p}&-J_q\end{pmatrix}\Delta,
 \qquad p+q=n,\quad p,q\ge1,
\end{equation}
where $\Delta$ is diagonal with entries $\pm1$.
Let $P$ and $Q$ be the sums of these switching signs in the two groups. Thus
$|P|\le p$, $|Q|\le q$, $P\equiv p\pmod2$, and $Q\equiv q\pmod2$.
Define
\[
 R_0=n-\frac{P^2}{p}-\frac{Q^2}{q},\qquad K=P^2+2PQ-Q^2.
\]
The nonzero spectrum of $A=cJ+rS$ is contained in the spectrum of
\begin{equation}\label{eq:compression}
 T=r\begin{pmatrix}p&\sqrt{pq}&0\\\sqrt{pq}&-q&0\\0&0&0\end{pmatrix}
      +cvv^T,
 \quad v=\begin{pmatrix}P/\sqrt p\\Q/\sqrt q\\\sqrt{R_0}\end{pmatrix}.
\end{equation}
A zero added in forming this compression can only increase the spread bound.
We also use \eqref{eq:compression} as a continuous family for real $P,Q$ in
the indicated rectangle. This relaxation is used to move to a particular
integer endpoint, not to discard the integer constraint.

For $c,r,R_0>0$ the eigenvalues have the form
$\lambda\ge\eta>0>-\mu$. The characteristic polynomial is
\begin{align}
 f(z)&=z^3-tz^2+e_2z+D,\label{eq:cubic}\\
 t&=cn+r(p-q),\nonumber\\
 e_2&=-2r^2pq+cr\{n(p-q)-K\},\qquad D=2cr^2pqR_0.\nonumber
\end{align}
Degenerate cases, including a possibly positive semidefinite rank-two limit,
are handled by continuity of the three-dimensional spread.

\section{An integer-preserving monotonicity reduction}
At fixed $R_0$, increasing $K$ decreases $e_2$ and increases both extreme
magnitudes:
\[
 \frac{\partial\lambda}{\partial K}
  =\frac{cr\lambda}{(\lambda+\mu)(\lambda-\eta)},\qquad
 \frac{\partial\mu}{\partial K}
  =\frac{cr\mu}{(\lambda+\mu)(\mu+\eta)}.
\]
As $e_2$ decreases the extreme roots move outwards and the middle positive
root moves downwards, so the relevant real-root branch does not collide.
Replacing $(P,Q)$ by $(|P|,|Q|)$ therefore cannot decrease spread. Assume below
that $P,Q\ge0$.

\begin{lemma}\label{lem:negative}
If $P+Q>0$, then $\mu\ge rqP/(P+Q)$.
\end{lemma}
\begin{proof}
Let $\gamma$ be minus the smallest eigenvalue of the sign part compressed
orthogonally to the all-ones vector. The $cJ$ term vanishes there, so
Rayleigh's principle gives $\mu\ge r\gamma$. For the continuous family
\eqref{eq:compression}, use $v^\perp$ in $\R^3$ in place of the all-ones
orthogonal complement; $\|v\|^2=n$, and the identical calculation applies.
The two potentially nonzero eigenvalues of this compression have trace
$p-q-K/n$ and product $-2pqR_0/n$. Thus $\gamma$ is the positive root of
\[
 h(g)=g^2+\left(p-q-\frac Kn\right)g-\frac{2pqR_0}{n}.
\]
Put $g_0=qP/(P+Q)$, $x=p-P$, $y=q-Q$, and $N=P+Q$. Direct expansion gives
$nN^2h(g_0)=-\mathcal R$, where
\begin{align*}
\mathcal R={}&P^2NQ^2
 +PQ(2P^2+5PQ+4Q^2)x
 +PQ(4P^2+9PQ+6Q^2)y\\
&+NQ(P+2Q)x^2+P(P^2+4PQ+5Q^2)y^2\\
&+2(P^3+4P^2Q+6PQ^2+2Q^3)xy\\
&+N(P+2Q)x^2y+(P^2+4PQ+2Q^2)xy^2+PQy^3.
\end{align*}
Every term is nonnegative. Hence $h(g_0)\le0$, $\gamma\ge g_0$, and the claim
follows. The identity is also checked by the supplied symbolic audit.
\end{proof}

\begin{lemma}\label{lem:monotone}
At fixed $p,q,Q\ge0$, the spread in \eqref{eq:compression} is nondecreasing
as $P$ runs from $0$ to $p$.
\end{lemma}
\begin{proof}
Write $s=\lambda+\mu$ and
\[
 C=2\lambda\mu+\eta(\lambda-\mu),\qquad W=\lambda-\mu-2\eta.
\]
Differentiating \eqref{eq:cubic} gives
\begin{equation}\label{eq:derivative}
 \frac{ds}{dP}=\frac{2cr\{C(P+Q)-2rqPW\}}
 {s(\lambda-\eta)(\mu+\eta)}.
\end{equation}
The denominator is positive. Also $C>0$: when $\lambda<\mu$, use
$\eta\le\lambda$ to see $C\ge\lambda\mu+\lambda\eta>0$.
If $W\le0$, the numerator is nonnegative. If $W>0$, then $\lambda>\mu$,
so Lemma~\ref{lem:negative} gives
\[
 C(P+Q)\ge2\lambda\mu(P+Q)\ge2\lambda rqP\ge2rqPW.
\]
The result follows; boundary points follow by continuity.
\end{proof}

We can therefore set $P=p$ while leaving $Q$ unchanged. Since $Q$ retains its
original parity and integrality, the numbers
\[
 y=\frac{q+Q}{2},\qquad z=\frac{q-Q}{2}
\]
are integers. The resulting endpoint matrix has base
\begin{equation}\label{eq:typeA}
 K_A(a)=\begin{pmatrix}1&1&a\\1&a&1\\a&1&a\end{pmatrix}
\end{equation}
and multiplicities $(p,y,z)$. When $z=0$ it is $B_q$. Otherwise $p,z\ge1$
and $y\ge1$. Its deficit has rank three when $p,y,z>0$; it is not a case of
the earlier rank-at-most-two deficit theorem.

\section{A three-type comparison for negative parameters}
For \eqref{eq:typeA}, let $q=y+z$, $n=p+q$, $b=1+a$, $d=1-a$, and $t=p+aq$.
The characteristic coefficients are
\[
 e_2=-dpq+dbz(p-q+z),\qquad e_3=-bd^2pyz.
\]
Write its eigenvalues as $\lambda\ge\eta\ge0>-\mu$ and $s=\lambda+\mu$.
An exact identity is
\begin{equation}\label{eq:etaidentity}
 s^2=F(q)-4dbz(p-q+z)+2t\eta-3\eta^2.
\end{equation}

\begin{lemma}\label{lem:eta}
For $-1\le a<1$, the middle eigenvalue satisfies $\eta\le dbz$.
\end{lemma}
\begin{proof}
Set $h_0=dbz$. The principal compression on the $p$ and $z$ groups is
\[
 R=\begin{pmatrix}p&a\sqrt{pz}\\a\sqrt{pz}&az\end{pmatrix},\quad
 g(w)=w^2-(p+az)w+adpz.
\]
Its smaller eigenvalue is at most $h_0$. For $a\le0$ this follows from its
nonpositive determinant. For $a\ge0$, a Rayleigh quotient orthogonal to
$(\sqrt p,\sqrt z)^T$ gives
$\lambda_{\min}(R)\le dpz/(p+z)\le dz\le dbz$.
Interlacing gives $\eta\le\lambda_{\max}(R)$, while direct substitution gives
\[
 f(h_0)=h_0\{g(h_0)-db^2yz\}.
\]
If $0<h_0<\eta$, then $g(h_0)\le0$ by its location between the eigenvalues
of $R$. This contradicts $f(h_0)>0$ between zero and the first positive root
of $f$. Degenerate cases follow by continuity.
\end{proof}

Suppose $-1\le a\le0$, so $0\le b\le1$. If $t\ge0$, define
$V=(1-b)F(q)+bF(y)$. Equation~\eqref{eq:etaidentity} gives
\[
 V-s^2=2t(dbz-\eta)+d^2bz^2+3\eta^2\ge0.
\]
Thus $s^2$ is bounded by a convex combination of candidate values. If $t<0$,
then $p<q$. With $\delta=q-p$, another exact identity gives
\[
 F(p)-s^2=db\left\{4\left(z-\frac\delta2\right)^2+2p\delta\right\}
           -2t\eta+3\eta^2\ge0.
\]
Consequently
\begin{equation}\label{eq:negativebound}
 s^2\le\max\{F(p),F(q),F(y)\},\qquad -1\le a\le0.
\end{equation}
For integer multiplicities these are candidate values. The harmless case
$y=0$ uses $F(0)=n^2\le M_n(a)^2$.

\section{An exact nonnegative-parameter certificate}
\begin{theorem}\label{thm:typeA}
For real $p,z\ge1$, $y\ge0$, and $0\le a<1$, the squared spread of the
three-dimensional compression of \eqref{eq:typeA} is at most
\[
 U=\frac{4(p+y+z)^2}{a+3}-\frac{(1-a)(a+3)}4.
\]
For integer multiplicities it is therefore at most $M_n(a)^2$.
\end{theorem}

\subsection{The spectral criterion}
Let $t,e_2,e_3$ be the characteristic coefficients of a real symmetric
$3\times3$ matrix. Put
\[
 D_0=t^2-3e_2,
\]
\[
 \mathcal D=t^2e_2^2-4e_2^3-4t^3e_3-27e_3^2+18te_2e_3.
\]
The three squared pairwise eigenvalue differences are the roots of
\[
 q(u)=u(u-D_0)^2-\mathcal D.
\]
Their largest value $s^2$ is at least $D_0$. Since
$q'(u)=(u-D_0)(3u-D_0)\ge0$ for $u\ge D_0$, it suffices to prove
\begin{equation}\label{eq:criterion}
 U\ge D_0,\qquad U(U-D_0)^2\ge\mathcal D.
\end{equation}
The relevant three-type coefficients are
\[
 t=p+a(y+z),\quad
 e_2=(1-a)\{-py+apz-(1+a)yz\},\quad
 e_3=-(1+a)(1-a)^2pyz.
\]

\subsection{Polynomial certificates and domain coverage}
Define the polynomials
\[
 B=\frac{4(a+3)(U-D_0)}{1-a},\qquad
 Q=\frac{64(a+3)^3\{U(U-D_0)^2-\mathcal D\}}{(1-a)^2}.
\]
All prefactors have positive sign on the parameter domain. Introduce
$a=T/(1+T)$ with $T\ge0$. The domain is covered by two charts:
\begin{center}
\begin{tabular}{llll}
\toprule
Chart & $p$ & $y$ & $z$\\
\midrule
Plus & $(a+2)v+u$ & $v$ & $1+w$\\
Minus & $1+v$ & $(1+v+u)/(a+2)$ & $1+w$\\
\bottomrule
\end{tabular}
\end{center}
Here $u,v,w\ge0$. The plus chart covers $p\ge(a+2)y$; the minus chart covers
$p\le(a+2)y$, $p\ge1$. After clearing the recorded positive denominators,
the plus-chart polynomials have positive coefficients. The minus-chart
polynomials are explicit sums of binomial squares with nonnegative monomial
multipliers and positive-coefficient remainders.

Every square entry $(\alpha,\beta,\gamma,k,w)$ represents
\[
 w\{4^kX^\alpha+4^{-k}X^\beta-2X^\gamma\},\qquad
 \alpha+\beta=2\gamma,\quad w>0.
\]
If $m=\min(\alpha,\beta)$ coordinatewise, this is exactly
\[
 wX^m\left(2^kX^{(\alpha-m)/2}-2^{-k}X^{(\beta-m)/2}\right)^2.
\]
All half-exponents are nonnegative integers and all coefficients are rational.
The positive-coefficient remainders are retained, not estimated numerically.
This proves \eqref{eq:criterion} and Theorem~\ref{thm:typeA}.

\subsection{Independent reconstruction}
The program \texttt{code/verify\_typeA\_certificate.py} uses only standard-library
integer and rational arithmetic. It does not import the generator or optimizer.
It independently reconstructs the spectral polynomials as follows. Let
$L=1+T$, $R=3+4T$, and express $p=p_N/m$, $y=y_N/m$, $z=z_N/m$ in either chart.
Define
\begin{align*}
 t_N&=Lp_N+T(y_N+z_N),\\
 e_{2,N}&=-Lp_Ny_N+Tp_Nz_N-(L+T)y_Nz_N,\\
 e_{3,N}&=-(L+T)p_Ny_Nz_N,\\
 D_N&=t_N^2-3e_{2,N},\\
 U_N&=16L^3(p_N+y_N+z_N)^2-m^2R^2,\\
 W_N&=U_N-4RD_N.
\end{align*}
Let $\mathcal D_N$ be the cubic discriminant evaluated at these numerator
coefficients. The independently reconstructed numerators are
\[
 Q_N=U_NW_N^2-64R^3\mathcal D_N,\qquad B_N=W_N.
\]
For the plus chart $m=L$, the certificate polynomials satisfy
$Q_N=L^3Q_+$ and $B_N=LB_+$. For the minus chart $m=2+3T$, they satisfy
$Q_N=m^2Q_-$ and $B_N=mB_-$. The verifier checks these identities before
checking every square and remainder coefficient. Thus it does not merely
trust that a supplied polynomial is the required spectral expression.

Floating-point linear programming was used only in searching for squares.
Its output was rationalized and independently checked exactly. The numerical
sweeps in the archive are supplementary audits, not premises of a proof.

\section{Completion of the centered-rank-two theorem}
Lemmas~\ref{lem:negative}--\ref{lem:monotone} reduce rank-two sign matrices
to \eqref{eq:typeA} with integer multiplicities. Equation~\eqref{eq:negativebound}
handles negative parameters; Theorem~\ref{thm:typeA} handles nonnegative ones.
If $S$ has rank one, write $S=\pm ss^T$ with $s_i\in\{-1,1\}$.
For the positive sign, $A=cJ+rss^T$ is positive semidefinite with trace $n$,
so its spread is at most $n$. For the negative sign its squared two-dimensional
spread is $n^2-4cr(\mathbf1^Ts)^2\le n^2$; degenerate cases also have spread
at most $n$. Since $F(1)>n^2$, neither case exceeds $M_n(a)$.
This proves Theorem~\ref{thm:centered}.

\section{Proof of the improved unrestricted strip}
Let $S$ be a symmetric sign matrix with rank at least three. Choose a nonsingular
$3\times3$ submatrix $T$. Such a sign matrix has singular values $2,2,1$.
To see this, its determinant is a nonzero multiple of four and is bounded in
absolute value by $\sqrt{27}$, so its absolute determinant is four. The Gram
matrix has diagonal entries three and off-diagonal entries $\pm1$; its
determinant forces eigenvalues $4,4,1$.

For any rank-at-most-two matrix $X$, the corresponding submatrix $R$ has a unit
null vector $v$. Hence
\[
 \|S-X\|_F^2\ge\|T-R\|_F^2\ge\|(T-R)v\|^2=\|Tv\|^2\ge1.
\]
Choose $X$ to retain the two extreme eigenvalues of $S$. The middle spectral
energy is at least one, whence
\[
 \spn(S)^2\le2(n^2-1).
\]
Write $L=\sqrt{2n^2-2}$ and $A=cJ+rS$. Subadditivity gives
$\spn(A)\le cn+rL$.
For $\epsilon=n\bmod2$ and $k=(n-\epsilon)/2$,
\[
 \spn(B_k)^2=n^2+r^2(n-\epsilon)^2+2r\epsilon(n-\epsilon).
\]
Put $C_\epsilon=(n-\epsilon)^2-(L-n)^2$. An exact comparison is
\begin{equation}\label{eq:stripidentity}
 \spn(B_k)^2-(cn+rL)^2=r\{(2-\epsilon)-cC_\epsilon\}.
\end{equation}
For even $n$, write $\beta=L/n$. Then
$C_0=n^2\beta(2-\beta)\le n^2$. For odd $n$,
\[
 C_1=2nL-2n^2-2n+3<2(\sqrt2-1)n^2.
\]
In both cases \eqref{eq:stripidentity} is nonnegative if
$c\le(1+\sqrt2)/(2n^2)$. Matrices with sign rank at most two are already
covered by Theorem~\ref{thm:centered}. Within this strip,
\[
 0\le\frac n2-\frac{n}{2+2c}
       =\frac{nc}{2(1+c)}\le\frac{1+\sqrt2}{4n}<\frac12,
\]
so $\lfloor n/2\rfloor$ is an allowed nearest integer in \eqref{eq:round}.
Thus the balanced candidate attains $M_n(a)$. Convexity extends the endpoint
comparison to every matrix in the entry box, proving Theorem~\ref{thm:strip}.

For $n\ge3$, the same calculation gives the sharper sufficient thresholds
\[
 c\le\begin{cases}
 \displaystyle\frac{2}{2n\sqrt{2n^2-2}-2n^2+2},& n\text{ even},\\[6pt]
 \displaystyle\frac{1}{2n\sqrt{2n^2-2}-2n^2-2n+3},& n\text{ odd}.
 \end{cases}
\]
They may be used with $a\le0$ so that the comparison remains entirely analytic.

\section{Relation to the supplied work and remaining gap}
The supplied earlier manuscripts treat continuous positive-semidefinite
deficits $J-A\succeq0$, binary deficits of rank at most two, and a narrower
boundary strip. They are retained verbatim under \texttt{provenance/}, not
silently rewritten as a complete solution. An inherited dimension-11
verification report is present, but the large certificate to which it refers
was not supplied. That finite enumeration is not represented as independently
reproduced by the present archive.

The new theorem concerns the rank of the \emph{centered sign matrix}
$S=(A-cJ)/r$, not the rank of the deficit $H=(J-A)/d$. These are distinct
conditions. Combining the present result with the supplied low-deficit-rank
results still leaves matrices for which both ranks are at least three and
the deficit is indefinite.

The missing universal inequality is
\[
 \boxed{\spn\bigl(J-(1-a)H\bigr)\le M_n(a)
 \quad\text{for every symmetric binary }H,
 \quad n\ge2,\quad -1\le a<1.}
\]
No argument here forces an unrestricted maximizer into the covered classes.
In particular, the rank-two Rayleigh difference $xx^T-yy^T$ does not force
its entrywise sign matrix to have rank two. Nested-neighborhood merging is
also not generally spread-increasing; the numerical exploration records
counterexamples to that proposed reduction, not counterexamples to SP-08.

\section{Reproducibility}
The archive includes this manuscript and PDF, Markdown proofs, exact
certificates, their independently checked rational identities, supplementary
numerical tests, search scripts, and inherited source files. The file
\texttt{STATUS.json} states the partial status explicitly. Start with
\begin{quote}\ttfamily
python code/verify\_typeA\_certificate.py
\end{quote}
which has no third-party dependencies. The symbolic audits use SymPy; the
numerical tests use NumPy. SciPy is needed only to rerun the certificate
search, not to validate its output. The machine-readable reports distinguish
exact verification from floating-point testing.

\begin{thebibliography}{9}
\bibitem{sp08} OpenProblemsInNLA contributors.
\emph{SP-08, eigenvalues and inverse problems}. Public problem statement:
\url{https://github.com/ajt60gaibb/OpenProblemsInNLA/tree/main/eigenvalues-and-inverse-problems/SP-08}.
\bibitem{supplied} Supplied research archive,
\emph{SP08\_partial\_results\_and\_proofs.zip}, and supplied manuscripts
\emph{sp08.tex}, \emph{boundary\_strip.tex}, and \emph{finite\_certificates.tex}.
Preserved in the archive's provenance directory. These are prior research
materials, not an assertion that unrestricted SP-08 is solved.
\end{thebibliography}
\end{document}
